
We now extend the semantics to several agents that run concurrently and interact through
shared memory and channels.
Two results carry over: replay soundness and crash recovery hold for whole systems, by a
reduction to the single-agent case (Section~\ref{sec:sys-replay}), and a probabilistic
noninterference theorem holds for concurrently scheduled agents under a floating-label
discipline (Section~\ref{sec:ni}).
Both depend on explicit restrictions, which we state as we go.

\subsection{Systems}
\label{sec:systems}

A \emph{system} consists of agents, shared memory, and channels.
An agent $\mathit{id}[M]_\Sigma$ has an identity, a term $M$, and a private state $\Sigma$ that
holds its local variables, its current label $\mathit{pc}(\Sigma)$
(Section~\ref{sec:ni}), a count of the oracle calls it has made, and a count of the agents it
has spawned.
Each agent runs a single thread; parallelism comes only from running several agents.
Shared memory $H$ has a fixed set of cells, each with a fixed label $\ell_k$; channels have
fixed labels $\ell_c$ and carry first-in-first-out queues.
Agent terms extend those of Section~\ref{sec:syntax} with
\[
  \mathbf{read}\;k\;\mathbf{as}\;x\,.\,M \;\mid\;
  \mathbf{write}\;k := e\,.\,M \;\mid\;
  \mathbf{send}\;c\;e\,.\,M \;\mid\;
  \mathbf{recv}\;c\;\mathbf{as}\;x\,.\,M \;\mid\;
  \mathbf{spawn}\;N\,.\,M .
\]
A spawned agent's identity is $(\mathit{id}, n)$, where $\mathit{id}$ is the parent and $n$ is
the parent's spawn count, so identities do not depend on any response.
A configuration is \emph{well formed} if its agents' identities are distinct and no agent has
an identity $(x, m)$ with $m$ at least the spawn count of $x$; steps preserve well-formedness,
so identities stay unique.
Every construct except an oracle call has a single successor.
A receive on an empty queue is disabled, and so is any step whose side condition
(Section~\ref{sec:ni}) fails: the agent blocks rather than aborting the system.
An agent whose term is $\mathbf{return}\;e$ halts.

A \emph{scheduler} $\sigma$ chooses, at each step, one enabled agent to step next, as a function
of the run so far, or stops the system when no agent is enabled.
We consider only \emph{deterministic} schedulers, and quantify over all of them.
For replay and recovery (Theorem~\ref{thm:sys}), a randomized scheduler is covered by making its
random choices oracle calls, so that they are drawn from a measure and journaled like any
other; the noninterference theorem of Section~\ref{sec:ni} requires a deterministic scheduler.

\begin{lemma}[Systems are programs]
\label{lem:sys-tree}
For every deterministic scheduler $\sigma$ and initial configuration $C_0$, the run of the
system is determined by the sequence of oracle responses: it is a tree in the sense of
Section~\ref{sec:syntax}, whose calls are all the agents' oracle calls in the order the
scheduler makes them.
\end{lemma}

\begin{proof}
The run so far is a function of the responses received so far, so at each step $\sigma$ picks
the agent as a function of those responses, and that agent's next construct either has a single
successor (pure computation, memory, channels, spawning) or is an oracle call.
Unfolding the run therefore gives a tree whose internal nodes are oracle calls and whose
deterministic steps are absorbed into the edges, as in the translation $T(M)$.
The tree is finite in depth: every step consumes a construct of some agent's term, and a spawn
splits $\mathbf{spawn}\;N.\,M$ into $N$ and $M$, so the total size of all agents' terms strictly
decreases.
A leaf is reached when no agent is enabled: either every agent has halted, or the remaining
agents are blocked; we label the leaf with the halted agents' results in the first case and
with a distinguished deadlock outcome in the second.
\end{proof}

\noindent
The Lean development formalizes this as a \emph{system}: a configuration type with a step
function that halts, takes a deterministic step, or makes an oracle call, unfolded into a tree
up to a step bound (\texttt{System.tree}).
It proves that raising the bound by one never changes a run that has already halted, provided
the halting output differs from the distinguished out-of-fuel output (\texttt{tree\_mono}).
For the calculus, a bound at least the total size of the initial terms covers the whole run.

\subsection{Replay and Recovery for Systems}
\label{sec:sys-replay}

Let all agents share a single journal that records every oracle call, whichever agent makes
it, in the order the calls occur.

\begin{theorem}[Replay soundness and crash recovery for systems]
\label{thm:sys}
For every deterministic scheduler and initial configuration, Theorem~\ref{thm:replay} and
Corollary~\ref{cor:resume} hold for the system's run with its global journal: replaying the
whole system from its initial configuration against the journal reproduces the run, and after
a crash that keeps any prefix of the journal, resuming the whole system with fresh randomness
gives the uninterrupted run on the spliced tape.
Since every step other than an oracle call is deterministic, the run reproduced includes all
memory, channel, and spawn steps, not only the oracle calls.
\end{theorem}

\begin{proof}
By Lemma~\ref{lem:sys-tree} the run is a tree, so Theorem~\ref{thm:replay} and
Corollary~\ref{cor:resume} apply to it directly; in Lean these are
\texttt{sys\_replay\_sound} and \texttt{sys\_crash\_resume}.
Given the responses, the deterministic steps between calls are the same in record, replay, and
resume, because the scheduler and every non-oracle construct are deterministic.
\end{proof}

\noindent
The theorem's restrictions are what make it simple, and each one rules out a failure that a
more ambitious design must handle.
\begin{itemize}[noitemsep]
  \item \emph{One journal, totally ordered.} A crash keeps a prefix of a single log, which is
    automatically a consistent state. With one journal per agent, persisted independently, a
    crash can keep an agent's receive without the matching send, and recovery must find a
    consistent cut of the logs~\citep{chandy1985snapshots} that also accounts for which memory
    write each read observed.
  \item \emph{Deterministic scheduling.} If replay could interleave agents differently from the
    recorded run, an agent could read a different value from shared memory than it did
    originally, and replay would diverge from the journal.
  \item \emph{Whole-system replay.} Replay re-runs every agent from the initial
    configuration. Restarting one agent in the middle of a run, as a supervisor might, would
    also require restoring that agent's local state and suppressing messages it had already
    sent; we do not treat it.
\end{itemize}

\subsection{Side Effects}
\label{sec:effects}

A side effect $\mathbf{do}\;f(e)$ changes the external world and returns a response.
We model it as two journaled oracle calls: an \emph{intent} call, whose measure is a point mass
on the record $(f, v)$ with $e \Downarrow v$, and a \emph{response} call, whose measure is the
world's response distribution for $f$ on $v$.
In record mode the runtime performs the effect between the two calls.
Replay reads both entries from the journal and performs nothing, so the replay part of
Theorem~\ref{thm:sys} covers effects: replay never repeats an effect.
The distributional part of crash recovery does not carry over, because the world's response to
an effect depends on the world's state, which is not a fixed measure per prompt.

Recovery is where effects differ from model calls.
A crash between the two entries leaves an intent without a response, and resuming draws the
response afresh, which in the runtime means performing the effect again; but the world may
already have applied it.
Whether this is safe is a property of the world, not of the journal.
If $f$ is idempotent on the world state and no other effect on the same state happened in
between, performing it again leaves the world as if it had been performed once; this is the
case the prototype handles by performing the effect again (Section~\ref{sec:eval-e2}).
Otherwise recovery needs a way to ask the world whether the effect happened, or a compensating
action with its own write-ahead record, and no policy that relies only on checking can
guarantee exactly-once effects without a bound on how long an effect may remain in
flight~\citep{li2026exactlyonce}.
We do not formalize the world, so these remain assumptions rather than theorems.

\subsection{Noninterference for Concurrent Agents}
\label{sec:ni}

Agents often handle data of different sensitivity: an agent that summarizes a confidential
document should not be able to leak it, even statistically, to an agent whose output is public.
We give a floating-label discipline and prove that, under a restricted class of schedulers, the
public behaviour of a system does not depend on its secret data.

\paragraph{Labels.}
Fix a lattice $(\mathcal{L}, \sqsubseteq)$ of security labels and a set $\Low \subsetneq
\mathcal{L}$ of observable labels that is downward closed and closed under joins.
Each agent carries a current label $\mathit{pc}$, which only rises.
The constructs of Section~\ref{sec:systems} have these side conditions and label updates:
\begin{center}
\small
\begin{tabular}{@{}lll@{}}
\toprule
\textbf{Construct} & \textbf{Permitted if} & \textbf{Then} \\
\midrule
$\mathbf{read}\;k$ & always & $\mathit{pc} := \mathit{pc} \sqcup \ell_k$ \\
$\mathbf{write}\;k := e$ & $\mathit{pc} \sqsubseteq \ell_k$ & cell $k$ keeps its label \\
$\mathbf{send}\;c\;e$ & $\mathit{pc} \sqsubseteq \ell_c$ & --- \\
$\mathbf{recv}\;c$ & $\mathit{pc} \sqsubseteq \ell_c$ & $\mathit{pc} := \ell_c$ \\
oracle call with output label $\ell$ & $\mathit{pc} \sqsubseteq \ell$ & $\mathit{pc} := \ell$ \\
$\mathbf{spawn}\;N$ & always & the child starts with the parent's $\mathit{pc}$ \\
\bottomrule
\end{tabular}
\end{center}
A receive is a destructive read, so it requires $\mathit{pc} \sqsubseteq \ell_c$ like a write.
Labels of cells and channels never change, and no cells or channels are created.
A step whose side condition fails is disabled, so the agent blocks
(Section~\ref{sec:systems}); aborting the whole system instead would let a secret decide whether
public agents run at all.
For this section we exclude side effects and clock reads, which expose the external world or
timing, and supervisors.
The global journal of Section~\ref{sec:sys-replay} records secret data, and its length reveals
how much secret computation happened, so it must be readable only at the top label.

\paragraph{Per-agent tapes.}
For noninterference it is convenient to give each agent its own tape: the response to the
$i$-th oracle call of agent $\mathit{id}$ with prompt $p$ is $\theta(\mathit{id}, i, p)$, for
a family $\theta$ whose coordinates are independent with $\theta(\mathit{id},i,p) \sim
\meas{p}$.
Because identities are unique and each agent's call count increases, distinct calls read
distinct coordinates, and the argument of Lemma~\ref{lem:tape} shows that this realizes
sampling from the oracle, for any deterministic scheduler.

\paragraph{Observations.}
A step is \emph{Low-visible} if the agent's current label before the step is in $\Low$ and the
data it touches is labelled in $\Low$: the cell for a read or write, the channel for a send or
receive, and the output label for an oracle call; spawning and halting are Low-visible when the
agent's label is in $\Low$.
Low-visible steps emit labels that carry their data: $\iota(k, v)$, $\omega(k, v)$, $c!v$,
$c?v$, $o(\mathit{id}, p, v)$, $\mathsf{spw}(\mathit{id}')$, and $\mathsf{done}(\mathit{id},
v)$.
The \emph{Low trace} $\mathit{tr}_{\Low}$ of a run is its sequence of Low-visible labels; other
steps are deleted, not replaced by a placeholder.
An agent is \emph{Low-active} if its label is in $\Low$ and its next construct is not a receive
on a channel whose label is outside $\Low$.
Two configurations are \emph{low-equivalent}, $C =_{\Low} C'$, if they have the same Low-active
agents, each with the same term and private state on both sides, the same values in cells
labelled in $\Low$, and the same queues on channels labelled in $\Low$.
Everything else, including agents that are not Low-active, may differ.

\paragraph{Assumptions.}
\begin{description}[noitemsep]
  \item[(A-SCHED)] The scheduler is deterministic, and at each step it chooses the
    highest-priority enabled agent under a total order on agent identities that is a function
    of the Low trace so far. It cannot observe steps that are not Low-visible, not even how
    many there were.
  \item[(A-TERM)] Every maximal run of consecutive steps by agents that are not Low-active is
    finite. In the calculus as defined every run is finite (Lemma~\ref{lem:sys-tree}), so this
    holds automatically; we state it because it is exactly what an extension with recursion
    would need, and it is the only termination assumption the proof uses.
  \item[(A-ORC)] Each oracle response depends only on the calling agent's prompt and its own
    tape coordinates. This holds in the calculus by construction; in a deployment it requires
    that model calls at different labels share no state, such as a cache or a conversation.
\end{description}

\begin{theorem}[Noninterference]
\label{thm:ni}
Assume (A-SCHED) and (A-ORC). For every tape family $\theta$ and all well-formed configurations
$C_0 =_{\Low} C_0'$, if (A-TERM) holds for the runs from $C_0$ and from $C_0'$ on $\theta$,
then the two runs have the same Low trace.
Consequently, for any distribution of tape families under which (A-TERM) holds almost surely
for both runs, the Low traces of the two runs have the same distribution; in particular this
holds for independent tapes, which realize sampling from the oracle.
\end{theorem}

\noindent
The proof rests on three lemmas.

\begin{lemma}[Steps that are not Low-active]
\label{lem:ni-high}
If $C =_{\Low} C'$ and $C$ steps to $D$ by an agent that is not Low-active, then
$D =_{\Low} C'$, the step is not Low-visible, and it changes no Low-active agent, no cell
labelled in $\Low$, and no channel labelled in $\Low$.
\end{lemma}

\begin{proof}
Suppose first that the agent's label $\mathit{pc}$ is not in $\Low$.
Because $\Low$ is downward closed, any label above $\mathit{pc}$ is outside $\Low$, so the side
conditions forbid writing a cell, sending or receiving on a channel, or making an oracle call
with an output label in $\Low$; reads only raise $\mathit{pc}$ further; and a spawned child
inherits $\mathit{pc}$ and is therefore not Low-active.
The step touches only the agent itself, cells and channels outside $\Low$, and possibly a new
agent that is not Low-active.
Otherwise the agent's label is in $\Low$ but its next construct is a receive on a channel
labelled outside $\Low$; the step changes only that channel and the agent, whose label becomes
$\ell_c \notin \Low$.
In both cases no Low-visible label is emitted and every component that $=_{\Low}$ constrains is
unchanged.
\end{proof}

\begin{lemma}[Low-active steps]
\label{lem:ni-low}
If $C =_{\Low} C'$ and the same Low-active agent takes its next step in both, then both steps
emit the same label, or neither is Low-visible, and the resulting configurations are
low-equivalent.
\end{lemma}

\begin{proof}
The agent has the same term and private state on both sides, so it performs the same construct
with the same values computed from its own state.
A read of a cell in $\Low$ returns the same value on both sides, and the agent's label stays in
$\Low$ because $\Low$ is closed under joins; a read of a cell outside $\Low$ may return
different values, but it raises the agent's label out of $\Low$ on both sides, so the agent is
no longer Low-active on either side.
A write or send stores the same value; if the target is in $\Low$ the label is the same, and
otherwise only components outside $\Low$ change.
A receive is on a channel in $\Low$ (the agent is Low-active), so it takes the same value from
the same queue.
An oracle call has the same prompt and is the agent's $i$-th call on both sides, so by
(A-ORC) it reads the same tape coordinate and gets the same response; if its output label is
outside $\Low$ the agent leaves the Low-active set on both sides.
A spawn creates the same identity with the same term and label on both sides.
Afterwards either the agent's term and state are the same on both sides, or its label is outside
$\Low$ on both sides; in either case it is Low-active on both sides or on neither, and the same
holds for a spawned child.
\end{proof}

\begin{lemma}[The next Low-active step]
\label{lem:ni-next}
Let $C =_{\Low} C'$ be reached by runs with the same Low trace, and suppose (A-TERM) holds.
Then either neither run takes another step by a Low-active agent, or both runs next take a
step by the same Low-active agent.
\end{lemma}

\begin{proof}
Whether a Low-active agent is enabled depends only on its term and private state (including its
label, which the side conditions check) and, for a receive, on a queue of a channel in $\Low$,
so the set $E$ of enabled Low-active agents is the same in $C$ and $C'$.
By Lemma~\ref{lem:ni-high}, steps by agents that are not Low-active do not change $E$, the
Low-active agents, or the Low trace, and the scheduler's priority order depends only on the Low
trace, so it is the same on both sides and stays fixed while such steps run.
If $E$ is empty, only agents that are not Low-active can step, and they can never enable a
Low-active agent, so neither run takes another Low-active step.
Otherwise let $X$ be the highest-priority agent in $E$.
On each side the scheduler fires the highest-priority enabled agent; until it fires $X$, it
can only fire agents above $X$ that are not Low-active, because no other Low-active agent above
$X$ is enabled.
By (A-TERM) such a sequence of steps is finite, and when it ends the highest-priority enabled
agent is $X$, which is still enabled.
So both runs next take a step by $X$.
\end{proof}

\begin{proof}[Proof of Theorem~\ref{thm:ni}]
Every Low-visible step is taken by a Low-active agent: an agent whose label is outside $\Low$
emits no Low-visible label, and a receive on a channel outside $\Low$ is not Low-visible.
By induction on the number of Low-active steps, using Lemmas~\ref{lem:ni-high}
to~\ref{lem:ni-next}, the two runs take the same sequence of Low-active steps with the same
labels, interleaved with steps that emit no Low-visible label, and remain low-equivalent with
the same Low trace after each.
Hence their Low traces are equal.
For the probabilistic statement, the two Low traces are equal on the event where (A-TERM) holds
for both runs; if this event has probability~1, they have the same distribution.
\end{proof}

\noindent
Each assumption rules out a known leak.
Without (A-SCHED), a scheduler that counts hidden steps can turn the running time of a secret
computation into a public difference~\citep{volpano1999probabilistic,sabelfeld2000probabilistic}.
Without (A-TERM), a secret computation that runs forever withholds every later public event,
the termination channel that termination-insensitive guarantees leave open.
Without fixed cell labels, an agent holding a secret could relabel a public cell and so decide,
through the secret, whether a public agent can read it; and with parallel threads inside an
agent, one thread's secret receive could block another thread's public output, which is why each
agent runs one thread.
Finally, the discipline is coarse: an agent that has read a secret can never again produce
public output, and the calculus has no construct for deliberate declassification.
LLMbda~\citep{garby2026llmbda} tracks labels on individual values for a single agent; combining
its finer tracking with concurrent scheduling is open.
